\section{Основные понятия и принципы математического моделирования. Прямые и обратные задачи математического моделирования. Универсальность математических моделей. Принцип аналогий. Иерархия моделей}

\subsection{Основные понятия математического моделирования}

\textbf{Модель} --- это объект, система или формальное описание, которое
замещает исследуемый объект и воспроизводит существенные для поставленной задачи
свойства этого объекта.

\textbf{Математическая модель} --- описание объекта или процесса с помощью
математических понятий и соотношений: уравнений, неравенств, функционалов,
алгоритмов, вероятностных законов и т.\,д.

В общем виде математическую модель можно представить как систему
\[
\mathcal{M} = (\mathbf{x}, \mathbf{p}, \mathbf{u},
\mathcal{F}, \mathcal{B}),
\]
где
\begin{itemize}
	\item $\mathbf{x}$ --- независимые переменные, например координаты и время;
	\item $\mathbf{u}$ --- искомые переменные, описывающие состояние системы;
	\item $\mathbf{p}$ --- параметры модели;
	\item $\mathcal{F}$ --- математические соотношения, описывающие поведение системы;
	\item $\mathcal{B}$ --- начальные, граничные и другие дополнительные условия.
\end{itemize}

Например, процесс теплопроводности может описываться уравнением
\[
\rho c\frac{\partial T}{\partial t}
=
\nabla\cdot(k\nabla T) + q,
\]
где $T$ --- температура, $\rho$ --- плотность, $c$ --- теплоёмкость,
$k$ --- коэффициент теплопроводности, $q$ --- объёмный источник тепла.

Математическое моделирование представляет собой исследование объекта не
непосредственно, а посредством построения и анализа его математической модели.

Основными целями моделирования являются:
\begin{itemize}
	\item объяснение поведения системы;
	\item прогнозирование её состояния;
	\item определение неизвестных параметров;
	\item управление системой;
	\item оптимизация её характеристик;
	\item проведение вычислительного эксперимента вместо дорогостоящего,
	опасного или невозможного натурного эксперимента.
\end{itemize}

При построении модели обычно используются следующие основные принципы.

\begin{enumerate}
	\item \textbf{Абстрагирование.}
	В модель включаются только свойства объекта, существенные с точки зрения
	поставленной задачи.
	
	\item \textbf{Адекватность.}
	Модель должна с требуемой точностью воспроизводить свойства реального объекта
	в рассматриваемом диапазоне условий.
	
	\item \textbf{Простота.}
	Не следует вводить в модель сложность, которая не улучшает решение
	поставленной задачи.
	
	\item \textbf{Системность.}
	Объект рассматривается как совокупность взаимосвязанных элементов.
	
	\item \textbf{Иерархичность.}
	Для одного объекта могут строиться модели различной сложности и детализации.
	
	\item \textbf{Проверяемость.}
	Результаты модели должны допускать сравнение с экспериментом,
	аналитическим решением или результатами других моделей.
\end{enumerate}


\subsection{Объект, модель и цель моделирования}

\textbf{Объект моделирования} --- реальная или абстрактная система,
процесс или явление, свойства которого требуется исследовать.

Перед построением модели необходимо определить \textbf{цель моделирования}.
От цели зависит, какие свойства объекта считаются существенными.

Например, автомобиль может моделироваться:
\begin{itemize}
	\item как материальная точка --- при исследовании движения по трассе;
	\item как твёрдое тело --- при исследовании динамики;
	\item как система упругих тел --- при расчёте прочности;
	\item как трёхмерное тело в потоке газа --- при исследовании аэродинамики.
\end{itemize}

Таким образом, одному объекту соответствует не одна ``истинная'' модель,
а множество моделей.

Удобно выделять три группы величин:
\begin{itemize}
	\item \textbf{входные воздействия} $\mathbf{f}$;
	\item \textbf{параметры модели} $\mathbf{p}$;
	\item \textbf{состояние системы} $\mathbf{u}$.
\end{itemize}

Тогда модель можно символически записать как
\[
\mathcal{A}(\mathbf{u},\mathbf{p}) = \mathbf{f}.
\]

Наблюдаемые величины могут быть связаны с состоянием системой наблюдения
\[
\mathbf{y} = \mathcal{H}(\mathbf{u}).
\]

Переход от реального объекта к математической модели всегда связан
с идеализацией. Например, могут использоваться предположения:
\begin{itemize}
	\item абсолютно твёрдое тело;
	\item несжимаемая жидкость;
	\item сплошная среда;
	\item идеальный газ;
	\item точечная масса;
	\item отсутствие трения;
	\item линейность зависимости между величинами.
\end{itemize}

Каждое такое предположение ограничивает область применимости модели.


\subsection{Этапы построения математической модели}

Типичная схема математического моделирования включает следующие этапы.

\begin{enumerate}
	\item \textbf{Постановка задачи.}
	
	Определяются объект, цель исследования, искомые величины и требуемая точность.
	
	\item \textbf{Выделение существенных свойств объекта.}
	
	Формулируются основные предположения и упрощения.
	
	\item \textbf{Построение концептуальной модели.}
	
	Определяются элементы системы и связи между ними.
	
	\item \textbf{Математическая формализация.}
	
	Законы и связи записываются в виде математических соотношений:
	\[
	\mathcal{A}(\mathbf{u},\mathbf{p})=\mathbf{f}.
	\]
	
	В физических задачах основой модели часто являются фундаментальные
	законы сохранения:
	\[
	\text{изменение запаса}
	=
	\text{приход}
	-
	\text{расход}
	+
	\text{источники}.
	\]
	
	\item \textbf{Задание начальных и граничных условий.}
	
	Одних дифференциальных уравнений обычно недостаточно для однозначного
	определения решения.
	
	\item \textbf{Определение параметров модели.}
	
	Они могут быть известны из эксперимента, справочных данных либо определяться
	в результате решения обратной задачи.
	
	\item \textbf{Исследование математической постановки.}
	
	Проверяются существование, единственность и устойчивость решения.
	
	\item \textbf{Выбор метода решения.}
	
	Если аналитическое решение невозможно, строится численный метод.
	
	\item \textbf{Реализация алгоритма в виде программы.}
	
	Таким образом возникает классическая триада
	\[
	\boxed{\text{модель} \longrightarrow
		\text{алгоритм} \longrightarrow
		\text{программа}}.
	\]
	
	\item \textbf{Верификация и валидация.}
	
	Проверяются правильность вычислительной реализации и соответствие модели
	реальному объекту.
	
	\item \textbf{Вычислительный эксперимент.}
	
	Выполняется серия расчётов при различных значениях параметров и условий.
	
	\item \textbf{Анализ результатов и уточнение модели.}
\end{enumerate}

Процесс моделирования обычно является циклическим:
\[
\text{объект}
\rightarrow
\text{модель}
\rightarrow
\text{расчёт}
\rightarrow
\text{сравнение с данными}
\rightarrow
\text{уточнение модели}.
\]


\subsection{Прямые задачи}

\textbf{Прямой задачей математического моделирования} называется задача,
в которой известны математическая модель, её параметры, начальные и граничные
условия, а требуется определить состояние системы.

В операторной форме:
\[
\mathcal{A}(\mathbf{u},\mathbf{p})=\mathbf{f},
\]
где известны $\mathcal{A}$, $\mathbf{p}$ и $\mathbf{f}$,
а требуется найти $\mathbf{u}$.

Например:
\begin{itemize}
	\item известны свойства материала и нагрузка --- требуется определить
	напряжения и деформации;
	\item известны коэффициенты теплопроводности и источники тепла ---
	требуется найти распределение температуры;
	\item известны начальное состояние и параметры динамической системы ---
	требуется определить её дальнейшую эволюцию.
\end{itemize}

Для задачи теплопроводности прямой задачей является, например,
\[
\begin{cases}
	\dfrac{\partial T}{\partial t}
	=
	a^2\dfrac{\partial^2 T}{\partial x^2},
	&
	0<x<L,\quad t>0,\\[2mm]
	T(x,0)=T_0(x),\\
	T(0,t)=T_1(t),\\
	T(L,t)=T_2(t),
\end{cases}
\]
где все коэффициенты и условия известны, а требуется определить
$T(x,t)$.

Классическим требованием к математической постановке является
\textbf{корректность по Адамару}. Задача называется корректной, если:
\begin{enumerate}
	\item решение существует;
	\item решение единственно;
	\item решение устойчиво относительно малых изменений исходных данных.
\end{enumerate}


\subsection{Обратные задачи}

В \textbf{обратной задаче} часть параметров, воздействий или характеристик
модели неизвестна и должна быть восстановлена по наблюдаемым результатам.

Пусть
\[
\mathcal{A}(\mathbf{u},\mathbf{p})=\mathbf{f},
\]
а измеряемая величина определяется оператором
\[
\mathbf{y}=\mathcal{H}(\mathbf{u}).
\]

В прямой задаче по заданному $\mathbf{p}$ находится $\mathbf{y}$:
\[
\mathbf{p}\longrightarrow\mathbf{u}
\longrightarrow\mathbf{y}.
\]

В обратной задаче по измеренным данным $\mathbf{y}^{\delta}$ требуется
восстановить
\[
\mathbf{p}.
\]

Например:
\begin{itemize}
	\item по измеренному температурному полю определить коэффициент
	теплопроводности;
	\item по перемещениям конструкции определить приложенные силы;
	\item по значениям поля на границе определить внутренний источник;
	\item по экспериментальным данным определить параметры модели;
	\item по томографическим измерениям восстановить внутреннюю структуру объекта.
\end{itemize}

Можно выделить несколько типов обратных задач:
\begin{itemize}
	\item восстановление параметров;
	\item восстановление начальных условий;
	\item восстановление граничных условий;
	\item восстановление источников;
	\item восстановление формы области;
	\item идентификация самой структуры модели.
\end{itemize}

Обратные задачи часто являются \textbf{некорректными}: малые ошибки измерений
могут вызывать большие изменения восстановленного решения.

Если необходимо решить
\[
A p = y^{\delta},
\]
где $y^{\delta}$ содержит погрешность измерения, непосредственное обращение
оператора $A$ может быть неустойчивым.

Поэтому применяют методы регуляризации. Например, в методе Тихонова
решается задача
\[
p_{\alpha}
=
\arg\min_p
\left(
\|Ap-y^{\delta}\|^2
+
\alpha\|Lp\|^2
\right),
\]
где $\alpha>0$ --- параметр регуляризации, а оператор $L$ задаёт
дополнительные требования к искомому решению.

Первое слагаемое обеспечивает согласование с экспериментальными данными,
а второе препятствует появлению неустойчивых решений.


\subsection{Универсальность моделей}

Одной из важнейших особенностей математического моделирования является
\textbf{универсальность математических моделей}.

Она состоит в том, что математически одинаковые или близкие модели могут
описывать совершенно различные физические, технические, биологические
или экономические процессы.

Причина заключается в том, что при абстрагировании от конкретной природы
объекта остаётся структура связей между величинами.

Например, уравнение
\[
\frac{\partial u}{\partial t}
=
D\Delta u
\]
может описывать:
\begin{itemize}
	\item теплопроводность;
	\item диффузию вещества;
	\item распространение примеси;
	\item некоторые процессы переноса в биологических системах.
\end{itemize}

При этом физический смысл $u$ и коэффициента $D$ различается,
но математическая структура задачи остаётся одинаковой.

Другой пример --- уравнение гармонического осциллятора
\[
\ddot{x}+\omega^2x=0.
\]

Оно возникает при описании:
\begin{itemize}
	\item механических колебаний;
	\item электрических колебаний;
	\item малых колебаний сложных динамических систем.
\end{itemize}

Универсальность позволяет:
\begin{itemize}
	\item переносить математические результаты между предметными областями;
	\item применять одни и те же численные методы к различным задачам;
	\item создавать универсальные программные комплексы;
	\item классифицировать модели не только по предметной области,
	но и по математической структуре.
\end{itemize}

Например, один и тот же метод конечных элементов может использоваться
для расчёта задач упругости, теплопроводности, электромагнетизма и других
краевых задач.


\subsection{Принцип аналогий}

\textbf{Принцип аналогий} состоит в том, что при построении модели нового объекта
используется сходство его структуры или поведения с уже изученной системой.

Различают:
\begin{itemize}
	\item физическую аналогию;
	\item структурную аналогию;
	\item математическую аналогию.
\end{itemize}

При математической аналогии процессы различной физической природы описываются
одинаковыми математическими уравнениями.

Например, стационарное распределение температуры без источников:
\[
\Delta T = 0
\]
и электростатический потенциал в области без зарядов:
\[
\Delta \varphi = 0
\]
описываются одним и тем же уравнением Лапласа.

Поэтому свойства решения, методы анализа и численные алгоритмы,
разработанные для одной задачи, могут быть использованы для другой.

Другой пример --- механическая система
\[
m\ddot{x}+c\dot{x}+kx=F(t)
\]
и электрическая $RLC$-цепь имеют математически аналогичные уравнения.

Принцип аналогий играет важную роль на начальном этапе моделирования:
вместо построения модели с нуля можно определить, к какому уже известному
классу математических моделей относится рассматриваемый объект.

Однако наличие математической аналогии не означает полной физической
эквивалентности систем. Область применимости модели должна проверяться
отдельно.


\subsection{Иерархия моделей}

Для одного объекта обычно существует не одна, а целая
\textbf{иерархия математических моделей}.

Под иерархией понимается совокупность моделей различной степени сложности,
детализации и вычислительной стоимости.

Например, движение тела можно описывать последовательно моделями:

\[
\boxed{\text{материальная точка}}
\rightarrow
\boxed{\text{абсолютно твёрдое тело}}
\rightarrow
\boxed{\text{упругое тело}}
\rightarrow
\boxed{\text{нелинейная модель материала}}.
\]

Каждый следующий уровень учитывает дополнительные эффекты.

Другой пример --- моделирование движения жидкости:
\[
\text{потенциальное течение}
\rightarrow
\text{уравнения Эйлера}
\rightarrow
\text{уравнения Навье--Стокса}
\rightarrow
\text{модели турбулентности}.
\]

Модели могут различаться:

\begin{itemize}
	\item по масштабу:
	\[
	\text{микро} \rightarrow \text{мезо} \rightarrow \text{макро};
	\]
	
	\item по пространственной детализации:
	\[
	\text{сосредоточенная}
	\rightarrow
	\text{одномерная}
	\rightarrow
	\text{двумерная}
	\rightarrow
	\text{трёхмерная};
	\]
	
	\item по типу описания:
	\[
	\text{детерминированная}
	\leftrightarrow
	\text{стохастическая};
	\]
	
	\item по характеру зависимостей:
	\[
	\text{линейная}
	\rightarrow
	\text{нелинейная};
	\]
	
	\item по учёту времени:
	\[
	\text{стационарная}
	\rightarrow
	\text{нестационарная}.
	\]
\end{itemize}

Выбор уровня иерархии определяется целью исследования.

Более сложная модель не обязательно лучше. Она требует:
\begin{itemize}
	\item большего количества исходных данных;
	\item большего количества параметров;
	\item более сложных численных методов;
	\item больших вычислительных ресурсов.
\end{itemize}

Поэтому модель должна иметь минимальную сложность, достаточную
для достижения требуемой точности.

На практике часто используется подход
\[
\text{простая модель}
\rightarrow
\text{сравнение с данными}
\rightarrow
\text{добавление необходимых эффектов}.
\]

Иерархия моделей также позволяет проводить взаимную проверку:
результаты сложной модели можно сравнивать с более простой моделью
в тех режимах, где применимы обе.


\subsection{Проверка и уточнение модели}

После построения математической модели необходимо убедиться,
что она пригодна для решения поставленной задачи.

Следует различать два понятия.

\textbf{Верификация} отвечает на вопрос:
\begin{quote}
	``Правильно ли мы решаем построенную математическую модель?''
\end{quote}

Проверяется:
\begin{itemize}
	\item правильность математических преобразований;
	\item правильность программной реализации;
	\item сходимость численного метода;
	\item зависимость результата от шага сетки;
	\item выполнение законов сохранения;
	\item совпадение с аналитическими решениями для тестовых задач.
\end{itemize}

Например, при последовательном измельчении сетки
\[
h_1 > h_2 > h_3 > \ldots
\]
численное решение должно стабилизироваться:
\[
u_{h} \rightarrow u,
\qquad h\rightarrow0.
\]

\textbf{Валидация} отвечает на другой вопрос:
\begin{quote}
	``Правильную ли модель реального объекта мы решаем?''
\end{quote}

Для этого результаты расчёта сравниваются с экспериментальными
или натурными данными.

Если наблюдения имеют вид $y_i^{\text{exp}}$, а модель даёт
$y_i^{\text{mod}}$, то, например, можно использовать среднеквадратичную ошибку
\[
E
=
\sqrt{
	\frac{1}{N}
	\sum_{i=1}^{N}
	\left(
	y_i^{\text{mod}}
	-
	y_i^{\text{exp}}
	\right)^2
}.
\]

Дополнительно проводится \textbf{анализ чувствительности}:
исследуется влияние параметров модели на результат,
\[
S_i
=
\frac{\partial y}{\partial p_i}.
\]

Если небольшое изменение некоторого параметра вызывает большое изменение
результата, точность определения этого параметра становится особенно важной.

Если результаты модели не соответствуют наблюдениям, возможны:
\begin{itemize}
	\item уточнение параметров;
	\item изменение начальных или граничных условий;
	\item отказ от слишком грубых предположений;
	\item добавление ранее не учитывавшихся физических эффектов;
	\item переход к более высокому уровню иерархии моделей.
\end{itemize}

Таким образом, математическое моделирование представляет собой
итерационный процесс:
\[
\boxed{
	\text{постановка задачи}
	\rightarrow
	\text{модель}
	\rightarrow
	\text{алгоритм}
	\rightarrow
	\text{расчёт}
	\rightarrow
	\text{проверка}
	\rightarrow
	\text{уточнение модели}
}.
\]

Основная идея состоит не в построении максимально сложной модели,
а в построении модели, которая при разумной вычислительной стоимости
адекватно описывает существенные для поставленной задачи свойства объекта.

\newpage